\section{Policy}\label{sec:method}
MorphMimic's controller is the policy $\pi_\theta(a_t\mid o_t,b)$ of
Section~\ref{sec:dataset}: it maps the present state, history, future reference
targets, and morphology to joint targets. This section defines its states and
actions, network, reward, and training.

\paragraph{State.}
Experiments use the \texttt{smpl\_mor} assets in Isaac Gym through
ProtoMotions~\cite{isaacgym,protomotions}, with simulation at 60\,Hz, two
substeps, and control decimation two, giving a 30\,Hz policy rate. The
observation combines the present state, history, future reference targets, a
historical action, and explicit morphology information. Let $B=24$ denote the
number of simulated bodies. The present state contains root height above
ground, root-relative body positions excluding the root, 6D body rotations, and
linear and angular velocities. Positions and velocities use a heading-aligned
frame and rotations remove the root heading, which removes dependence on
arbitrary global heading and horizontal origin while retaining height and body
configuration. We add seven historical states at control-step offsets
$\{1,2,3,4,8,16,32\}$, covering about 1.07\,s and each in its own root/heading
frame, and four future reference slots at $\{1,2,4,8\}$, covering about
0.27\,s. Future slots contain target positions relative to the current root,
heading-aligned target rotations, position differences from the current body
state, relative rotations $\smash{R_{t,j}^{-1}\hat{R}_{t+k,j}}$ (expressed with
respect to the current body rotation), and heading-aligned linear and angular
velocity differences. The history exposes recent changes and the future targets
a short anticipation window; contact flags, an explicit phase scalar, text, and
estimated physical feature vectors are not part of the observation.
Table~\ref{tab:observations} summarizes the inputs of actor and critic, with
dimensions per slot, positions in metres, linear velocities in m/s, angular
velocities in rad/s, and dimensionless 6D rotations.

\begin{table}[t]
\centering
\small
\caption{Observation inputs supplied to both actor and critic.}\label{tab:observations}
\begin{tabularx}{\linewidth}{@{}lrlY@{}}
\toprule
Input & Dimension & Slots & Representation and path \\
\midrule
Current state & $15B-2=358$ & 1 &
Height; $3(B-1)$ positions; $6B$ rotations; $3B$ linear and $3B$ angular
velocities. Current-state encoder. \\
History & $358$ & 7 &
Same state representation at the seven past offsets. Shared history encoder. \\
Future reference & $24B=576$ & 4 &
Target and relative poses ($18B$), plus velocity differences ($6B$).
Shared future encoder. \\
Historical action & $D=69$ & 1 &
A raw previous policy-action vector, before the tanh and joint-target mapping.
Bypasses attention. \\
Morphology & $11$ & 1 &
$[g,\beta_1,\ldots,\beta_{10}]$ from the assigned asset.
Bypasses attention. \\
\bottomrule
\end{tabularx}
\end{table}

\paragraph{Action.}
The policy outputs a raw action $a_t\in\R^D$ with $D=69$, mapped to joint
targets by
\begin{equation}
    q_t^{\mathrm{target}}
      = q^{\mathrm{offset}}
        + q^{\mathrm{scale}}\odot\tanh(a_t),
    \label{eq:action_mapping}
\end{equation}
with offsets and scales derived from the configured joint ranges. The targets
drive the simulator's built-in position proportional--derivative (PD)
controller; the action is not a residual added to the reference pose. PD
stiffness, damping, and effort limits are scaled per joint by the mass of the
associated rigid body relative to that body in the reference asset (base values
in Appendix~\ref{app:control}). This body-level mass scaling is not a single
scale based on total character mass, and it does not guarantee dynamically
equivalent actuation across geometrically different bodies.

\paragraph{Network.}
MorphMimic uses temporal slot/type attention. Three observation encoders, one
for the current state, one shared across history slots, and one shared across
future slots, are multilayer perceptrons (MLPs) with two 256-unit hidden layers,
rectified linear unit (ReLU) activations, and a 256-dimensional token output, and their inputs use running mean/standard-deviation
normalization with a magnitude clamp of five. The resulting 12 tokens (one
current, seven historical, four future) enter a Transformer as
\begin{equation}
    z_i = E_{\tau(i)}(x_i)
        + e_i^{\mathrm{slot}} + e_{\tau(i)}^{\mathrm{type}},
    \qquad i=0,\ldots,11,
    \label{eq:token_embedding}
\end{equation}
where $\tau(i)$ identifies the current, history, or future source: slot
embeddings distinguish the configured temporal offsets, and type embeddings the
observation sources. The Transformer has two layers, four attention heads, token
width 256, and feed-forward width 1,024~\cite{attention}, and its current-state
token supplies the pooled output after the configured output projection and
activation. Figure~\ref{fig:network} traces the encoded observations through
attention and the separate actor and critic heads, which share input categories
but not parameters.

\begin{figure}[t]
\centering
\fbox{\begin{minipage}{0.94\linewidth}
\small
\begin{tabularx}{\linewidth}{@{}YcY@{}}
Current state: $1\times358$ & $\longrightarrow$ &
Current encoder: $1\times256$ token \\
Past states: $7\times358$ & $\longrightarrow$ &
History encoder: $7\times256$ tokens \\
Future targets: $4\times576$ & $\longrightarrow$ &
Future encoder: $4\times256$ tokens
\end{tabularx}
\begin{center}
$\downarrow$\\
12 tokens + learned temporal-slot and source-type embeddings\\
$\downarrow$\\
2 Transformer layers, 4 heads, width 256\\
$\downarrow$\\
Pool current token: $h_t\in\R^{256}$\\
$\downarrow$\\
Concatenate $[h_t,\ a^{\mathrm{hist}}_t,\ g,\boldsymbol{\beta}]$
and normalize
\end{center}
\begin{tabularx}{\linewidth}{@{}YY@{}}
\textbf{Actor:} six 2,896-unit layers
$\longrightarrow$ 69 action means
$\longrightarrow$ tanh/range scaling
$\longrightarrow$ PD targets. &
\textbf{Critic:} four 1,024-unit layers
$\longrightarrow$ scalar state value.
\end{tabularx}
\smallskip
The entire encoder/attention pathway is instantiated separately for actor
and critic; the drawing summarizes the common design, not shared weights.
Morphology and action bypass attention into each head. Text is absent.
\end{minipage}}
\caption{Temporal slot/type controller and separate actor/critic heads.}\label{fig:network}
\end{figure}

The final head input concatenates the 256-dimensional attention output, the
historical raw action, and the 11-dimensional morphology vector; this
336-dimensional input also uses running normalization and clamp five. Betas are
not manually divided by three, embedded as a separate attention token, or used
for adaLN modulation in the selected model. Actor and critic have separate token
encoders, attention modules, slot/type embeddings, and MLP heads, with six
hidden layers of width 2,896 for the actor and four of width 1,024 for the
critic, all with ReLU activations; they output a Gaussian action mean and a
scalar value, respectively. PPO samples actions with fixed diagonal log standard
deviation $-2.9$, and evaluation uses the mean action.

\paragraph{Reward.}
The training reward combines global body tracking and root-height tracking. For
corresponding simulated and reference body positions $p_j,\hat{p}_j$, rotations
$R_j,\hat{R}_j$, and velocities $v_j,\hat{v}_j$ and $\omega_j,\hat{\omega}_j$,
define
\begin{align}
    e_p &= \frac{1}{3B}\sum_{j=1}^{B}\|p_j-\hat{p}_j\|_2^2,
    \\
    e_R &= \frac{1}{B}\sum_{j=1}^{B}
                  d_{\SO(3)}(R_j,\hat{R}_j)^2, \\
    e_v &= \frac{1}{3B}\sum_{j=1}^{B}\|v_j-\hat{v}_j\|_2^2,
    \\
    e_\omega &= \frac{1}{3B}\sum_{j=1}^{B}
                     \|\omega_j-\hat{\omega}_j\|_2^2, \\
    e_h &= (h-\hat{h})^2,
\end{align}
where $d_{\SO(3)}$ is the geodesic rotation angle in radians and the
coordinate-wise means in the position and velocity errors match the implemented
reductions. The reward is
\begin{equation}
\begin{split}
    r_t ={}& 0.5\exp(-25e_p) + 0.3\exp(-5e_R)
             + 0.1\exp(-0.5e_v)\\
           &+0.2\exp(-0.1e_\omega) + 0.2\exp(-100e_h),
\end{split}\label{eq:reward}
\end{equation}
with weights that are not normalized to sum to one. There is no explicit
effort, action-smoothness, or contact-matching reward in this experiment, and no
source-dependent switch between canonical and HUMOS reward functions.

\paragraph{Optimization.}
MorphMimic is trained directly across all bodies on refined HUMOS references,
with no neutral pretraining stage, using PPO~\cite{ppo}. References are sampled
only from variants matching the environment's asset.
Reference-state initialization starts at frame zero with probability 0.2 and
otherwise at a random valid reference time. Training episodes last at most
1,000 control steps and terminate on falls (a 0.15\,m height setting, with
permitted foot-contact bodies), while tracking-error termination is disabled, so
this rule is distinct from the evaluation failure threshold of
Section~\ref{sec:metrics}. PPO uses 32-step rollouts, discount 0.99, generalized
advantage estimation (GAE) parameter 0.95, policy-ratio clipping 0.2, normalized
advantages, and separate Adam optimizers with learning rates $2\times10^{-5}$ for
the actor and $10^{-4}$ for the critic. Other settings and the actual run budgets
are in Appendix~\ref{app:training}.

\paragraph{Streamed motion sampling.}
The full-scale experiment streams per-clip files through
\texttt{GlobalClipPool}. The training manifest is deterministically partitioned
across GPU ranks, and each rank keeps a bounded resident pool containing all 128
variants of its selected clips, with environment-level motion sampling
respecting the assigned morphology. The pool is rebuilt every 256 epochs: most
slots are selected by priority, combining estimated clip difficulty and an
exploration bonus, and a fraction 0.2 is reserved for uniform random selection
to maintain rehearsal outside the highest-priority set. This is a training
curriculum, not uniform sampling of the complete corpus. Validation membership
is fixed independently of pool rebuilding and does not update sampling weights,
so resident-pool and validation metrics measure different populations. Sampling
constants and the pool size actually used are in
Appendices~\ref{app:reproducibility} and~\ref{app:history}.
